\documentclass[12pt]{article}
\usepackage[utf8]{inputenc}
\usepackage[T1]{fontenc}
\usepackage{hyperref}
\title{Tsimshian Language Resources for K-12 Curriculum Development}
\author{Compiled by Claude Sonnet 5 (claude-sonnet-5) for Andr\'as Kornai}
\date{}
\begin{document}
\maketitle

\section{Introduction}
Tsimshian (ISO 639-3 tsi) is documented here through its Coast Tsimshian variety, Sm'algyax. One conference report in the search results says that of about 10,000 Tsimshian people, no more than roughly 175 are fluent. The sources I found are mainly dictionaries and talking-dictionary tools. I did not find a current school curriculum.

\section{Primary Educational Resources}
\begin{itemize}
\item Sm'algyax Student Dictionary, described on the UBC RelLex wiki: \url{https://wiki.ubc.ca/Documentation:RelLex/Sm%27algyax_Student_Dictionary}
\end{itemize}
No textbooks or graded primers were found in the results. The results mention new school curriculum sessions at a conference, with no details.

\section{Linguistic and Reference Materials}
\begin{itemize}
\item Sm'algyax: A Reference Dictionary and Grammar for the Coast Tsimshian Language (John A. Dunn, University of Washington Press and Sealaska Heritage Foundation, 1995; the two parts are separate volumes). The UBC knowledge base entry: \url{https://knowledgebase.arts.ubc.ca/?p=2593}
\item UBC RelLex wiki entry for the same work: \url{https://wiki.ubc.ca/Documentation:RelLex/Sm%27algyax:_A_Reference_Dictionary_and_Grammar_of_the_Coast_Tsimshian_Language}
\item Dictionary of Shm'algyack, compiled by Donna May Roberts with Elders of Metlakatla, Alaska, posted free by Sealaska Heritage Institute as a searchable PDF (news report): \url{https://ictnews.org/archive/native-language-dictionary-added-sealaskas-free-offerings/}
\end{itemize}

\section{Religious and Community Resources}
I did not find scripture or community-organisation pages in these results.

\section{Digital and Audio Resources}
\begin{itemize}
\item Sm'algyax Living Legacy Talking Dictionary, an online multimedia dictionary (UBC knowledge base): \url{https://knowledgebase.arts.ubc.ca/?p=2589}
\item UBC RelLex wiki entry for the Talking Dictionary: \url{https://wiki.ubc.ca/Documentation:RelLex/Sm%E2%80%99algya%CC%B1x_Living_Legacy_Talking_Dictionary}
\item Sm'algyax on FirstVoices (UBC RelLex wiki entry). The search summary says an app exists on Apple and Google stores: \url{https://wiki.ubc.ca/Documentation:RelLex/Sm%27algya%CC%B1x_FirstVoices}
\end{itemize}

\section{Curriculum Development Context}
There is a standard reference dictionary, and Dunn's spelling is reported to be the authority for teachers. There are also open talking-dictionary tools. The documented gaps are graded readers, textbooks and a published curriculum. The results mention School District 52 (Prince Rupert) as a source of earlier dictionary drafts. Its present curriculum should be investigated directly.

\end{document}